\documentclass[spanish,openany]{report}
% \usepackage{darkmode} \enabledarkmode
\usepackage{wrapfig}
\usepackage{enumitem}
\input{~/git/Clase/template/preamble.tex}
\input{~/git/Clase/template/macros.tex}
\input{~/git/Clase/template/letterfonts.tex}
\graphicspath{{/home/kuma/git/Clase/Fp_Gs_Robotica/images/}}
\setimagefolder{/home/kuma/git/Clase/Fp_Gs_Robotica/images/}
\date{\today}
% \hbadness=99999
\usepackage{microtype}
% \addbibresource{./bibliografia.bib}
\usepackage{titling}
\usepackage{afterpage}


% Header and footer
% \fancypagestyle{plain}{
%     \fancyhf{} %Clear Everything.
%     \fancyfoot[C]{Página \thepage\ de \pageref{LastPage} }%Page Number
%     \renewcommand{\headrulewidth}{1pt} %0pt for no rule, 2pt thicker etc...
%     \renewcommand{\footrulewidth}{1pt}
%     \fancyfoot[R]{\small{Daniel Carbajales Soto}}
%     \fancyhead[R]{\small{Cinta Transportadora con 3 Velocidades}}
%     \fancyhead[L]{\small{\today}}
% }
% \pagestyle{plain}

\begin{document}

%----------------------------------------------------------------------------------------
%	                            PORTADA
%----------------------------------------------------------------------------------------
% \makeportada{}{Daniel Carbajales Soto}{ 1° CFGS Automatización y Robótica Industrial}{}

%----------------------------------------------------------------------------------------
%                                    HEADER 
%----------------------------------------------------------------------------------------
\header{SEHN - Practica 1 - Cuadro Trifásico}{Daniel Carbajales Soto}{1° CFGSARI}
%----------------------------------------------------------------------------------------
%                               TABLE OF CONTENTS
%----------------------------------------------------------------------------------------
% \maketoc
%----------------------------------------------------------------------------------------
%	                            TABLES
%----------------------------------------------------------------------------------------
% this is just a demo in case i need to make a table
%
%
%
% \begin{tabularx}{\linewidth}{|C{2cm}|Y|Y|C{2cm}|}
% \hline
% \makecell{Header 1} & \makecell{Header 2 \\ with linebreak} & \makecell{Header 3} & \makecell{Header 4} \\
% \hline
% Row 1 & This is a long cell text that automatically wraps in the Y column & Another long cell & Short \\
% \hline
% Row 2 & \multirow{2}{*}{Multirow text} & More text & Short \\
% \cline{1-1} \cline{3-4}
% Row 3 & & Extra text & Short \\
% \hline
% \end{tabularx}


%----------------------------------------------------------------------------------------
%	                            DOCUMENT
%----------------------------------------------------------------------------------------


\section*{ejemplo de representación}
Imagina que debes realizar el montaje de un cuadro de distribución para tres líneas monofásicas de 230 V, a partir de una red de 
distribución trifásica con neutro de 400 V. La línea general debe estar protegida mediante un interruptor magnetotérmico y un
interruptor diferencial. Además, cada una de las líneas monofásicas debe disponer de interruptores magnetotérmicos individuales. 
Dibuja las representaciones unifilar y multifilar.
\section*{Monataje del cuadro trifásico}
\begin{figure}[h]
    \centering
    \includegraphics[width=0.7\textwidth]{cuadro_trifasico.png}
    \caption{Monataje del cuadro trifásico}
    \label{fig:cuadro_trifasico}
\end{figure}

\clearpage
	
\section*{Representación unifilar}
\begin{figure}[h]
    \centering
    \incfigpdf[0.45]{SEHN_prac_1_esq_uni}
    \caption{Representación unifilar}
    \label{fig:cuadro_trifasico_unifilar}
\end{figure}
\clearpage
\section*{Representación multifilar}
\begin{figure}[h]
    \centering
    \incfigpdf[0.45]{SEHN_prac_1_esq_mult}
    \caption{Representación multifilar}
    \label{fig:cuadro_trifasico_multifilar}
\end{figure}


%%%%%%%%%%%%%%%%%%%%%%%%%%%%%%%%%%%%%%%%%%%%%%%%%%%%%%%%%%
%%%%%%%%%%%%%%%%%%%% bibliography %%%%%%%%%%%%%%%%%%%%%%%%
% \nocite{*}                                 %%%%%%%%%%%%%
% \printbibliography                         %%%%%%%%%%%%%
%%%%%%%%%%%%%%%%%%%%%%%%%%%%%%%%%%%%%%%%%%%%%%%%%%%%%%%%%%
%%%%%% Last page in case there is no bibliography %%%%%%%%
% \label{Last Page}                           %%%%%%%%%%%%%
%%%%%%%%%%%%%%%%%%%%%%%%%%%%%%%%%%%%%%%%%%%%%%%%%%%%%%%%%%

\end{document}
